% An error tex RECOVERS from prints three lines: the message, the line the
% mouth had reached, and a second line aligned to where in it the scanner
% stopped (tex.web §82's error, §311's show_context). The two lines are
% trimmed to fit -- line one keeps its TAIL, because the interesting end is
% where the scan stopped, and line two keeps its HEAD, each marking a cut with
% `...' (§317, half_error_line=50 and error_line=79). Both trims are here.
%
% The position is per DIGIT, not per constant: §445 tests before folding a
% digit in, so the split falls after the ten nines that fit and before the
% eleventh, which is the one that overflowed. A radix constant overflows by
% the same rule with no slack at all.
%
% Recovery is `error' and carry on: the register keeps 2147483647, the rest of
% the line is read normally, and the run ends with the line tex writes when
% something was reported.
\catcode`\{=1 \catcode`\}=2
\count1=1 \count2=2 \count3=3 \count4=4 \count5=5 \count6=99999999999 \count7=7 \count8=8 \count9=9 \count10=10 \count11=11 \count12=12 \count13=13
\count14="FFFFFFFF
\message{[\the\count6][\the\count13][\the\count14]}
\end
